\chapter{同一个G条件，现在有三个不同的临界问题}\label{ch:targets}
\section{允许的分母从一个素数变成一组素数}
本篇统一用 $k$ 表示每次平均的元数，$n$ 表示总位置数。非恒定有理输入照前篇中心化、本原化，得到整数零和 $X$，记差分最大公因数为 $G$。

共同余数障碍的证明不变，只需把“$\ell\ne p$”改为“$\ell\nmid k$”。因此一般必要条件是
\begin{equation}\label{eq:G}
G\mid n,\qquad \rad(G)\mid k.
\end{equation}
$\rad(G)$ 表示 $G$ 的不同素因子的乘积。第二项就是“$G$ 的每个素因子都整除 $k$”，并不要求 $G\mid k$。

例如六平均允许 $G=2^a3^b$，十平均允许 $G=2^a5^b$。同一个条件在不同位置数上会化简：六平均十元有 $G\mid10$，所以允许 $G=1,2$；六平均九元有 $G\mid9$，所以可能的 $G=1,3,9$ 全部被允许。这两个例子不是用了不同猜想。

\section{自动满足条件，不等于已经证明可解}
把命题“在 $n$ 个位置上，所有有理输入都恰由\eqref{eq:G}判定”记为 $\mathcal P(k,n)$。若 $\rad(n)\mid k$，由于 $G\mid n$，条件对每个输入自动成立。此时 $\mathcal P(k,n)$ 断言的就是\emph{所有输入都能平均}，仍然是一项需要证明的结论。

最直接的警示是24平均27元。27的素因子只有三，三整除24，所以G条件对所有输入都自动合法。然而
\[((-1)^{24},0,0,24)
\]
不能归零。下一讲会证明这个反例。可见“分母允许”远远不等于“位置上的操作足够”。

\section{M、H、N分别问哪件事}
排除一步全体平均的平凡维数 $n=k$，定义\src{S01}
\begin{align*}
M(k)&=\min\{n>k:\mathcal P(k,n)\},\\
H(k)&=\min\{n>k:\mathcal P(k,n),\ \rad(n)\nmid k\},\\
N(k)&=\min\{N>k:\text{每个 }n\ge N\text{ 都满足 }\mathcal P(k,n)\}.
\end{align*}
$M$ 问第一次在哪个位置数成功；$H$ 只看G条件真正能筛掉输入的维数；$N$ 则要求以后不再失败。

二平均有 $M=N=4$，但四元G自动合法，所以 $H=5$。六平均已有 $M=9,H=10$；仍未证明 $N=9$ 或 $N=10$。十平均同样有 $M=16,H=17$，最后阈值还没闭合。

不要从两个连续成功点推出后面全成功。若在零和配置中只平均一个子系统，它的均值一般不是零；冻结外部坐标再调用小维数解法，可能会到错误的目标。前篇的整体共同尺度约定在这里依然重要。

\section{合数真正改变的三件事}
第一，非零差分模 $k$ 未必可逆。例如模六的二不是单位。原来用 $i(a-b)$ 任意调余数的步骤失效，但这个公因子也可能提供新的整数混合比例。

第二，最早候选位置数可能远小于 $2k$。九平均的候选临界是十三，放不下两个九重块。平方族要改用 $(k,\sqrt k,1)$，不能沿用素数端点的 $(k,k,1)$。

第三，多个允许素数的分母变化必须由\emph{同一条实际路径}协调。CRT能拼接同余条件，却不能替一张超库存的表增加位置，也不能把两条各自有效的局部路径自动合并。

本篇后面的三项未完成任务正对应这些困难：较大素数奇次幂的临界核心、混合素因子的临界点、以及已解点之后的全部维数接续。\src{S02}\src{S03}
